\documentclass[10pt,aspectratio=169]{beamer}
\usepackage[utf8]{inputenc}
\usepackage[T1]{fontenc}
\usepackage[french]{babel}
\usepackage{amsmath}
\usepackage{amssymb}
\usepackage{booktabs}
\usepackage{array}
\usepackage{listings}
\usepackage{tikz}
\usetikzlibrary{positioning,calc,arrows.meta}
\usepackage{eurosym}

% Thème : celui des autres séances
\usetheme{Madrid}
\usecolortheme{default}

\definecolor{bleuactu}{RGB}{0,51,102}
\definecolor{vertactu}{RGB}{0,102,51}
\definecolor{rougeactu}{RGB}{153,0,0}
\definecolor{orangeactu}{RGB}{230,126,34}
\definecolor{codegreen}{rgb}{0,0.6,0}
\definecolor{codegray}{rgb}{0.5,0.5,0.5}
\definecolor{codepurple}{rgb}{0.58,0,0.82}
\definecolor{backcolour}{rgb}{0.95,0.95,0.92}

\setbeamercolor{structure}{fg=bleuactu}
\setbeamercolor{frametitle}{bg=bleuactu,fg=white}

\lstset{
    basicstyle=\ttfamily\scriptsize,
    breaklines=true,
    frame=single,
    backgroundcolor=\color{backcolour},
    commentstyle=\color{codegreen},
    keywordstyle=\color{blue},
    stringstyle=\color{codepurple},
    language=Python,
    showstringspaces=false
}

\newcommand{\fx}[1]{\texttt{\small #1}}

% Compilation pdfLaTeX sans erreur : caractères absents de l'encodage T1, accents dans les listings
\DeclareUnicodeCharacter{2248}{\ensuremath{\approx}}
\DeclareUnicodeCharacter{2264}{\ensuremath{\leq}}
\DeclareUnicodeCharacter{2265}{\ensuremath{\geq}}
\lstset{literate=
  {é}{{\'e}}1 {è}{{\`e}}1 {ê}{{\^e}}1 {à}{{\`a}}1 {â}{{\^a}}1 {î}{{\^i}}1 {ô}{{\^o}}1 {ù}{{\`u}}1
  {û}{{\^u}}1 {ç}{{\c c}}1 {É}{{\'E}}1 {€}{{\texteuro}}1 {≥}{{$\geq$}}1 {≤}{{$\leq$}}1
  {«}{{\guillemotleft}}1 {»}{{\guillemotright}}1}

\title[Séance 3 -- Excel, Python et agents d'IA]{Séance 3 : Excel, Python et agents d'IA}
\subtitle{Les outils de l'actuaire, appliqués au contrôle prudentiel}
\author[F. Borel-Mathurin]{Fabrice BOREL-MATHURIN}
\institute[ACPR-DEAR et SAF Lyon 1]{
    ACPR - DEAR\\
    et\\
    Laboratoire SAF Université Claude Bernard Lyon 1
}
\date{2026-2027}

\begin{document}

\frame{\titlepage}

\begin{frame}{Plan de la séance (4 h)}
\tableofcontents
\end{frame}

\begin{frame}{Objectifs pédagogiques}
\begin{block}{À l'issue de cette séance, vous serez capable de :}
\begin{enumerate}
  \item Organiser un classeur Excel qu'un autre peut auditer, et y agréger des SCR avec une matrice de corrélation
  \item Écrire en Python un contrôle reproductible sur des données prudentielles, et le rapprocher du tableur
  \item Confier une tâche à un agent d'IA : la cadrer, puis relire et vérifier son résultat
  \item Citer les règles qui encadrent ces usages : qualité des données, responsabilité, confidentialité, règlement sur l'IA, DORA
\end{enumerate}
\end{block}
\vspace{0.3cm}
\begin{alertblock}{Déroulé}
\textbf{Cours :} 2 h (Excel 30 min, Python 40 min, agents 40 min) \quad | \quad
\textbf{TD 3 :} 2 h, un projet en trois parties
\end{alertblock}
\end{frame}

\begin{frame}{Rappel de la séance 2}
\begin{block}{Trois questions, sans notes}
\begin{enumerate}
  \item Pourquoi un écart d'acquisition vaut-il zéro au bilan prudentiel ?
  \item D'où vient un impôt différé actif, et à quelle condition est-il retenu ?
  \item Pourquoi un SCR non couvert n'entraîne-t-il pas immédiatement un retrait d'agrément ?
\end{enumerate}
\end{block}
\vspace{0.2cm}
\small Réponses en 5 minutes, à deux, puis mise en commun.
\end{frame}

% =====================================================================
\section{Pourquoi les outils comptent en prudentiel}
% =====================================================================

\begin{frame}{Des chiffres sous contrainte}
\begin{columns}[T]
\begin{column}{0.48\textwidth}
\begin{block}{Ce qu'on demande au chiffre prudentiel}
\begin{itemize}
  \item \textbf{Volumes} : la liste des actifs (S.06.02) compte une ligne par titre détenu
  \item \textbf{Délais} : remises annuelles et trimestrielles à date fixe
  \item \textbf{Qualité} : des données « appropriées, exhaustives et exactes » pour les provisions techniques (directive, art.~82)
  \item \textbf{Traçabilité} : l'auditeur, la fonction actuarielle et l'ACPR doivent pouvoir refaire le calcul
\end{itemize}
\end{block}
\end{column}
\begin{column}{0.48\textwidth}
\begin{alertblock}{Quand l'outil se trompe}
\small
\begin{itemize}
  \item \textbf{JPMorgan, 2012} (« London Whale ») : le modèle de VaR reposait sur des tableurs
        alimentés à la main ; une formule divisait par une somme au lieu d'une moyenne et
        sous-estimait le risque (rapport interne, 2013)
  \item \textbf{Reinhart et Rogoff, 2010} : une plage de cellules incomplète dans un tableur a
        biaisé un résultat de politique économique très cité (Herndon, Ash et Pollin, 2013)
\end{itemize}
\end{alertblock}
\end{column}
\end{columns}
\end{frame}

\begin{frame}{Trois outils, trois usages}
\begin{table}
\centering
\small
\renewcommand{\arraystretch}{1.5}
\begin{tabular}{p{2.4cm}p{4.6cm}p{5.6cm}}
\toprule
\textbf{Outil} & \textbf{Pour} & \textbf{Limites} \\
\midrule
\textbf{Tableur} & Explorer, prototyper, montrer un calcul pas à pas & Erreurs invisibles, pas d'historique, difficile à tester, 1\,048\,576 lignes au plus \\
\textbf{Python} & Automatiser, traiter du volume, tester, versionner & Demande d'apprendre ; un script faux se trompe vite et partout \\
\textbf{Agent d'IA} & Écrire du code, des tests, une première version & Se trompe avec assurance ; ne remplace ni la vérification ni la signature \\
\bottomrule
\end{tabular}
\end{table}
\begin{center}
\small Le bon réflexe : \textbf{deux calculs indépendants qui se rapprochent}, par exemple le tableur et le script.
\end{center}
\end{frame}

% =====================================================================
\section{Excel express}
% =====================================================================

\begin{frame}{Un classeur qu'un autre peut auditer}
\begin{center}
\begin{tikzpicture}[
  onglet/.style={draw=bleuactu, fill=bleuactu!8, rounded corners, minimum width=2.6cm,
                 minimum height=1cm, align=center, font=\small},
  fleche/.style={-{Latex[length=2mm]}, thick, bleuactu}]
\node[onglet] (e) {\textbf{Entrées}\\\scriptsize données, paramètres};
\node[onglet, right=0.8cm of e] (c) {\textbf{Calculs}\\\scriptsize une formule par colonne};
\node[onglet, right=0.8cm of c] (s) {\textbf{Sorties}\\\scriptsize ce qu'on remet};
\node[onglet, right=0.8cm of s, draw=rougeactu, fill=rougeactu!8] (k) {\textbf{Contrôles}\\\scriptsize VRAI = anomalie};
\draw[fleche] (e) -- (c); \draw[fleche] (c) -- (s); \draw[fleche] (s) -- (k);
\end{tikzpicture}
\end{center}
\begin{itemize}
  \item \textbf{Aucune constante en dur} : 15\,\%, 50\,\%, 0,5 point vont dans des cellules nommées (\emph{Formules > Gestionnaire de noms})
  \item \textbf{Mettre sous forme de tableau} (Ctrl + L) : références lisibles, formules étendues automatiquement
  \item \textbf{Une seule formule par colonne}, recopiée sur toutes les lignes
  \item \textbf{Repérer les antécédents} (\emph{Formules > Audit de formules}) avant de faire confiance à une cellule
  \item \textbf{Des contrôles visibles} : totaux qui se recoupent, bilan équilibré, mise en forme conditionnelle
\end{itemize}
\end{frame}

\begin{frame}{Les formules qui servent}
\begin{table}
\centering
\small
\renewcommand{\arraystretch}{1.3}
\begin{tabular}{lll}
\toprule
\textbf{Français} & \textbf{Anglais} & \textbf{Usage prudentiel} \\
\midrule
\fx{SOMMEPROD} & \fx{SUMPRODUCT} & Somme de flux actualisés, de montants pondérés \\
\fx{RECHERCHEX} & \fx{XLOOKUP} & Retrouver un taux de la courbe, un facteur par segment \\
\fx{SI}, \fx{ET}, \fx{OU} & \fx{IF}, \fx{AND}, \fx{OR} & Écrire un contrôle \\
\fx{MIN}, \fx{MAX} & \fx{MIN}, \fx{MAX} & Appliquer une limite d'éligibilité, un corridor \\
\fx{NB.SI} & \fx{COUNTIF} & Compter les anomalies \\
\fx{PRODUITMAT}, \fx{TRANSPOSE} & \fx{MMULT}, \fx{TRANSPOSE} & Agréger des SCR avec une matrice de corrélation \\
\fx{LET} & \fx{LET} & Nommer les étapes d'une formule longue \\
\bottomrule
\end{tabular}
\end{table}
\small En français, les arguments sont séparés par un point-virgule et les décimales par une virgule.
\end{frame}

\begin{frame}{Exemple : limites d'éligibilité et agrégation du SCR}
\begin{block}{Fonds propres éligibles à la couverture du SCR (règlement délégué, art.~82)}
\small
\fx{t3\_elig = MIN([@tier3]; Lim\_T3*[@scr])} \hfill niveau 3 < 15\,\% du SCR\\
\fx{t2\_elig = MIN([@tier2]; Lim\_T2T3*[@scr]-[@t3\_elig])} \hfill niveaux 2 + 3 ≤ 50\,\%\\
\fx{ratio = ([@tier1]+[@t2\_elig]+[@t3\_elig])/[@scr]}
\end{block}
\begin{block}{SCR de base : $\sqrt{s^\top C\, s}$ (directive, annexe IV)}
\small
\fx{=RACINE(PRODUITMAT(PRODUITMAT(TRANSPOSE(s);C);s))}

\smallskip
Marché 60, contrepartie 5, vie 30 : BSCR $= 75{,}17$, contre 95 en simple somme. La
diversification vaut 19,83.
\end{block}
\end{frame}

% =====================================================================
\section{Python express}
% =====================================================================

\begin{frame}{Pourquoi passer au code}
\begin{columns}[T]
\begin{column}{0.48\textwidth}
\begin{block}{Ce que le code apporte}
\begin{itemize}
  \item \textbf{Reproductible} : relancé, il redonne le même chiffre
  \item \textbf{Lisible} : chaque étape est écrite, nommée, commentée
  \item \textbf{Testable} : des assertions vérifient les résultats connus
  \item \textbf{Versionné} : git garde chaque modification, son auteur et sa raison
  \item \textbf{Sans limite de lignes} pratique
\end{itemize}
\end{block}
\end{column}
\begin{column}{0.48\textwidth}
\begin{block}{Où l'écrire}
\begin{itemize}
  \item \textbf{Google Colab} : un notebook dans le navigateur, rien à installer
  \item \textbf{Sur votre machine} : Python 3.10 ou plus, puis \fx{pip install numpy pandas matplotlib jupyter}
  \item \textbf{Notebook} pour explorer et expliquer ; \textbf{script} pour la production
\end{itemize}
\end{block}
\end{column}
\end{columns}
\vspace{0.2cm}
\begin{center}\small Support : notebook \emph{Python express}, sur la page de la séance (« Ouvrir dans Colab »).\end{center}
\end{frame}

\begin{frame}[fragile]{Calculer et vectoriser}
\begin{columns}[T]
\begin{column}{0.48\textwidth}
\small Une fonction, comme une formule qu'on n'écrit qu'une fois :
\begin{lstlisting}
def valeur_actuelle(flux, taux, date):
    """Valeur actuelle d'un flux."""
    return flux / (1 + taux) ** date

valeur_actuelle(100, 0.03, 5)   # 86.26
\end{lstlisting}
\end{column}
\begin{column}{0.48\textwidth}
\small NumPy calcule sur des vecteurs entiers :
\begin{lstlisting}
import numpy as np
flux  = np.array([40, 40, 40, 1040])
dates = np.arange(1, 5)
taux  = np.array([.025, .028, .030, .032])
(flux / (1 + taux) ** dates).sum()

s = np.array([60, 5, 30, 0, 0])
np.sqrt(s @ C @ s)              # 75.17
\end{lstlisting}
\end{column}
\end{columns}
\end{frame}

\begin{frame}[fragile]{Manipuler un tableau avec pandas}
\begin{lstlisting}
import pandas as pd
d = pd.read_csv("bilans-assureurs.csv", sep=";")          # une feuille, des colonnes nommées

d["excedent"] = (d["actif"] - d["provisions_techniques"]
                 - d["dettes_subordonnees"] - d["autres_passifs"])   # une colonne calculée
d[d["excedent"] != d["excedent_declare"]]                  # filtrer les écarts
d[d.isna().any(axis=1)]                                    # chercher les données manquantes

t3 = np.minimum(d["tier3"], 0.15 * d["scr"])               # une limite sur toute la colonne
d.to_excel("controles.xlsx", index=False)                  # rendre le résultat au tableur
\end{lstlisting}
\small Pas de boucle sur les lignes : chaque opération s'applique à la colonne entière.
\end{frame}

\begin{frame}[fragile]{Vérifier}
\begin{columns}[T]
\begin{column}{0.48\textwidth}
\begin{block}{Assertions}
\begin{lstlisting}
assert np.isclose(bscr, 75.17, atol=0.01), \
       "BSCR différent du calcul Excel"
\end{lstlisting}
\small Si le résultat change, le script s'arrête : une régression se voit tout de suite.
\end{block}
\end{column}
\begin{column}{0.48\textwidth}
\begin{block}{Les pièges classiques}
\small
\begin{itemize}
  \item \fx{0.1 + 0.2 == 0.3} est faux : comparer avec une \textbf{tolérance}
  \item \fx{dropna()} fait disparaître des lignes \textbf{en silence}
  \item Un chemin absolu (\fx{C:/Users/…}) ne marche que sur une machine
  \item \fx{> 1} au lieu de \fx{>= 1} : un ratio de 100\,\% devient « non couvert »
\end{itemize}
\end{block}
\end{column}
\end{columns}
\vspace{0.2cm}
\begin{alertblock}{Le rapprochement}
Le script et le tableur, construits séparément, doivent donner les mêmes chiffres, au centime près.
\end{alertblock}
\end{frame}

\begin{frame}{Versionner et vérifier automatiquement}
\begin{itemize}
  \item \textbf{git} garde l'historique : qui a changé quoi, quand, pourquoi
  \item \textbf{Une branche par tâche, une demande de fusion (\emph{pull request})} : une seconde personne relit avant que le changement n'entre
  \item \textbf{L'intégration continue} relance les vérifications à chaque modification : compilation, exécution des notebooks, tests
\end{itemize}
\vspace{0.3cm}
\begin{block}{Exemple : le site de ce cours}
Chaque diapositive, notebook et corrigé est dans un dépôt git. À chaque modification, les PDF sont
recompilés, chaque notebook est exécuté, les liens sont vérifiés, puis le site est publié. Un
notebook qui ne tourne plus est signalé dès la modification qui l'a cassé.
\end{block}
\end{frame}

% =====================================================================
\section{Piloter un projet avec des agents d'IA}
% =====================================================================

\begin{frame}{Assistant, agent : de quoi parle-t-on ?}
\begin{columns}[T]
\begin{column}{0.48\textwidth}
\begin{block}{Un assistant}
Répond à une question dans une conversation. Vous copiez, collez, exécutez.
\end{block}
\begin{block}{Un agent}
Reçoit un objectif, découpe le travail, \textbf{agit} avec des outils (lire et écrire des
fichiers, exécuter du code, chercher), observe le résultat et recommence jusqu'à ce qu'il
estime avoir fini. Exemples : Claude Code, GitHub Copilot en mode agent.
\end{block}
\end{column}
\begin{column}{0.48\textwidth}
\begin{alertblock}{Ce qu'il faut savoir}
\begin{itemize}
  \item Il produit vite du travail plausible, \textbf{y compris quand il se trompe}
  \item Il ne connaît que ce que vous lui donnez et ce qu'il trouve
  \item Il comble les manques de la consigne par des hypothèses
  \item Il peut affirmer avoir vérifié ce qu'il n'a pas vérifié
\end{itemize}
\end{alertblock}
\end{column}
\end{columns}
\end{frame}

\begin{frame}{Le cycle d'un projet mené avec un agent}
\begin{center}
\begin{tikzpicture}[
  etape/.style={draw=bleuactu, fill=bleuactu!8, rounded corners, minimum width=2.7cm,
                minimum height=1.1cm, align=center, font=\small},
  fleche/.style={-{Latex[length=2mm]}, thick, bleuactu}]
\node[etape] (a) {\textbf{1. Cadrer}\\\scriptsize objectif, critères};
\node[etape, right=1.1cm of a] (b) {\textbf{2. Déléguer}\\\scriptsize une tâche courte};
\node[etape, right=1.1cm of b] (c) {\textbf{3. Vérifier}\\\scriptsize exécuter, rapprocher};
\node[etape, right=1.1cm of c] (d) {\textbf{4. Documenter}\\\scriptsize tracer, faire relire};
\draw[fleche] (a) -- (b); \draw[fleche] (b) -- (c); \draw[fleche] (c) -- (d);
\draw[fleche] (c.south) -- ++(0,-0.6) -| node[pos=0.25, below, font=\scriptsize]{corriger la consigne, recommencer} (b.south);
\end{tikzpicture}
\end{center}
\vspace{0.2cm}
\begin{itemize}
  \item \textbf{Découper} : des tâches courtes, chacune vérifiable seule (« ajouter le contrôle C4 », pas « refaire le reporting »)
  \item \textbf{Garder la main} sur ce qui engage : interprétation des textes, hypothèses, validation, signature
  \item \textbf{Déléguer} ce qui se vérifie : code, tests, graphiques, mise en forme, premier jet d'un document
\end{itemize}
\end{frame}

\begin{frame}[fragile]{Une bonne consigne}
\begin{lstlisting}[language={},basicstyle=\ttfamily\scriptsize]
OBJECTIF   Ajouter au script controles.py le contrôle C4 : « MCR hors de 25-45 % du SCR »
           (directive 2009/138/CE, art. 129).
CONTEXTE   Données : bilans-assureurs.csv, fictives, en M€, séparateur « ; ».
CONTRAINTES pandas et NumPy seulement ; ne pas modifier les autres contrôles ;
           ne supprimer aucune ligne ; comparer les nombres avec une tolérance.
ACCEPTATION C4 signale l'entité D et elle seule ; une assertion le vérifie ;
           le script s'exécute depuis le dossier de la séance.
NE PAS     citer un article sans l'avoir vérifié ; affirmer que tout est détecté.
LIVRER     les modifications, la sortie du script, la liste des vérifications faites.
\end{lstlisting}
\small Les critères d'acceptation disent \textbf{à quoi vous reconnaîtrez un travail juste}. Sans
eux, l'agent décide seul quand il a fini.
\end{frame}

\begin{frame}{Vérifier ce que l'agent a produit}
\begin{columns}[T]
\begin{column}{0.48\textwidth}
\begin{block}{La liste de relecture}
\begin{enumerate}
  \item L'exécuter soi-même, sur une autre machine
  \item Rapprocher d'un calcul indépendant (le tableur)
  \item Tester les cas limites : donnée manquante, valeur exactement au seuil
  \item Vérifier \textbf{chaque} référence réglementaire dans le texte
  \item Lire les modifications ligne à ligne
  \item Se méfier des conclusions générales (« tout est vérifié »)
\end{enumerate}
\end{block}
\end{column}
\begin{column}{0.48\textwidth}
\begin{alertblock}{Retour d'expérience}
\small
Relecture d'un support de cours préparé en 2025 avec des assistants d'IA :
\begin{itemize}
  \item matrice de corrélation du SCR inexacte
  \item ajustement symétrique sans son terme de $-8\,\%$, et de sens inversé
  \item références réglementaires introuvables
  \item marqueurs d'outil oubliés rendant un notebook illisible
  \item chemins de fichiers d'une autre machine
\end{itemize}
Tout était plausible. Seuls les textes et l'exécution du code l'ont révélé.
\end{alertblock}
\end{column}
\end{columns}
\end{frame}

\begin{frame}{Les règles du jeu dans une entreprise d'assurance}
\begin{table}
\centering
\small
\renewcommand{\arraystretch}{1.35}
\begin{tabular}{p{3cm}p{10cm}}
\toprule
\textbf{Question} & \textbf{Repère} \\
\midrule
Quelles données ? & Seulement ce que la politique de l'entreprise autorise pour l'outil ; données personnelles : RGPD (règlement (UE) 2016/679) ; rien de confidentiel dans un service non approuvé \\
Qui répond du chiffre ? & L'entreprise et la personne qui valide. Même logique que l'externalisation : l'entreprise reste pleinement responsable (directive, art.~49) \\
Quelle gouvernance ? & Un système de gouvernance garantissant une gestion saine et prudente (art.~41) ; revue par une seconde personne ; données de qualité (art.~82) \\
Quels textes propres à l'IA ? & Règlement (UE) 2024/1689 sur l'IA, applicable par étapes : l'évaluation des risques et la tarification en assurance vie et santé sont « à haut risque » (annexe~III) \\
Quel risque informatique ? & DORA, règlement (UE) 2022/2554, applicable depuis le 17 janvier 2025 : un service d'IA en nuage est un prestataire tiers de services informatiques \\
\bottomrule
\end{tabular}
\end{table}
\end{frame}

% =====================================================================
\section{TD 3 et synthèse}
% =====================================================================

\begin{frame}{TD 3 : contrôler les états de six assureurs}
\begin{columns}[T]
\begin{column}{0.32\textwidth}
\begin{block}{A. Tableur (40 min)}
\small
Recalculer excédent, fonds propres éligibles, ratios ; sept contrôles ; agréger un SCR.
\end{block}
\end{column}
\begin{column}{0.32\textwidth}
\begin{block}{B. Python (40 min)}
\small
Automatiser les mêmes contrôles, traiter les cas limites, rapprocher du tableur.
\end{block}
\end{column}
\begin{column}{0.32\textwidth}
\begin{block}{C. Agent (40 min)}
\small
Relire un script produit par une IA ; rédiger la consigne d'un agent ; règles du jeu.
\end{block}
\end{column}
\end{columns}
\vspace{0.4cm}
\small Données : \fx{bilans-assureurs.csv} et \fx{matrice-correlation.csv} sur la page de la séance.
À rendre : le classeur, le script, la consigne et une note de revue d'une page.
\begin{alertblock}{Avant de commencer}
Les six ratios déclarés sont-ils tous justes ? Pariez, puis vérifiez.
\end{alertblock}
\end{frame}

\begin{frame}{À retenir}
\begin{enumerate}
  \item Un chiffre prudentiel doit pouvoir être \textbf{refait} : entrées séparées, paramètres nommés, contrôles visibles.
  \item Le tableur sert à comprendre, le code à \textbf{reproduire et tester} ; les deux se rapprochent au centime.
  \item Un agent d'IA va vite, se trompe avec assurance, et \textbf{ne vérifie pas à votre place}.
  \item Une bonne consigne contient ses \textbf{critères d'acceptation} ; une bonne revue exécute, rapproche et vérifie chaque référence.
  \item L'entreprise et la personne qui valide restent \textbf{responsables} ; données, gouvernance, IA et risque informatique ont leurs textes.
\end{enumerate}
\end{frame}

\begin{frame}{Pour la prochaine séance}
\begin{block}{À faire}
\begin{itemize}
  \item Terminer le TD 3 ; le corrigé est publié sur le site après la séance
  \item Revoir les unités 4 à 6 du \textbf{polycopié de prérequis} : flux aléatoire, flux optionnel, contrat d'assurance
\end{itemize}
\end{block}
\begin{block}{Questions de rappel, en ouverture de la séance 4}
\begin{itemize}
  \item Pourquoi une limite d'éligibilité peut-elle faire passer un ratio sous 100\,\% ?
  \item Citez deux pièges qui font échouer un script en silence.
  \item Que doit contenir la consigne d'un agent pour qu'on puisse juger son travail ?
\end{itemize}
\end{block}
\end{frame}

\begin{frame}[c]
\centering
\Huge Questions ?

\vspace{1cm}
\Large
\textbf{Séance 4} : meilleure estimation non-vie
\end{frame}

\end{document}
